\chapterbackmatter{Weinberg Exercises}

\begin{enumerate}
\WeinbergExercise{1} Show that if the zero-momentum coefficient functions satisfy the
conditions \eqref{eq:5.1.23} and \eqref{eq:5.1.24}, then the coefficient functions \eqref{eq:5.1.21}
and \eqref{eq:5.1.22} for arbitrary momentum satisfy the defining conditions
Eqs.~\eqref{eq:5.1.19} and \eqref{eq:5.1.20}.

\WeinbergExercise{2} Consider a free field $\psi_\ell^\mu(x)$ which annihilates and creates a
self-charge-conjugate particle of spin $\tfrac32$ and mass $m\ne0$. Show how
to calculate the coefficient functions $u_\ell^\mu(\mathbf p,\sigma)$, which
multiply the annihilation operators $a_{\mathbf p,\sigma}$ in this field, in
such a way that the field transforms under Lorentz transformations like a
Dirac field $\psi_\ell$ with an extra four-vector index $\mu$. What field
equations and algebraic and reality conditions does this field satisfy?
Evaluate the matrix $P_{\ell m}^{\mu\nu}(p)$, defined (for $p^2=-m^2$) by
\[
  \sum_\sigma u_\ell^\mu(\mathbf p,\sigma)
    u_m^{\nu *}(\mathbf p,\sigma)
  \equiv (2p^0)^{-1}P_{\ell m}^{\mu\nu}(p).
\]
What are the commutation relations of this field? How does the field
transform under the inversions $P,C,T$?

\WeinbergExercise{3} Consider a free field $h^{\mu\nu}(x)$ satisfying
$h^{\mu\nu}(x)=h^{\nu\mu}(x)$ and $h^\mu{}_\mu(x)=0$, which annihilates and
creates a particle of spin two and mass $m\ne0$. Show how to calculate the
coefficient functions $u^{\mu\nu}(\mathbf p,\sigma)$, which multiply the
annihilation operators $a_{\mathbf p,\sigma}$ in this field, in such a way that
the field transforms under Lorentz transformations like a tensor. What field
equations does this field satisfy? Evaluate the function
$P^{\mu\nu,\kappa\lambda}(p)$, defined by
\[
  \sum_\sigma u^{\mu\nu}(\mathbf p,\sigma)
    u^{\kappa\lambda *}(\mathbf p,\sigma)
  \equiv (2p^0)^{-1}P^{\mu\nu,\kappa\lambda}(p).
\]
What are the commutation relations of this field? How does the field
transform under the inversions $P,C,T$?

\WeinbergExercise{4} Show that the fields for a massless particle of spin $j$ of type
$(A,A+j)$ or $(B+j,B)$ are the $2A$th or $2B$th derivatives of fields of
type $(0,j)$ or $(j,0)$, respectively.

\WeinbergExercise{5} Work out the transformation properties of fields of transformation
type $(j,0)+(0,j)$ for massless particles of helicity $\pm j$ under the
inversions $P,C,T$.

\WeinbergExercise{6} Consider a generalized Dirac field $\psi$ that transforms according to
the $(j,0)+(0,j)$ representation of the homogeneous Lorentz group. List the
tensors that can be formed from products of the components of $\psi$ and
$\psi^\dagger$. Check your result against what we found for $j=\tfrac12$.

\WeinbergExercise{7} Consider a general field $\psi_{ab}$ describing particles of spin $j$
and mass $m\ne0$, that transforms according to the $(A,B)$ representation of
the homogeneous Lorentz group. Suppose it has an interaction Hamiltonian of
the form
\[
  H_{\mathrm{int}}=\int d^3x\,
  \bigl[\psi_{ab}(x)J^{ab}(x)
    +J^{ab\dagger}(x)\psi_{ab}^\dagger(x)\bigr],
\]
where $J^{ab}$ is an external c-number current. What is the asymptotic
behavior of the matrix element for emitting these particles for energy
$E\gg m$ and definite helicity? (Assume that the Fourier transform of the
current has values for different $a,b$ that are of the same order of
magnitude, and that do not depend strongly on $E$.)
\end{enumerate}
